\begin{theorem}[Support-to-nerve realization]
\label{thm:philosophy-support-to-nerve-realization}
The assignment
$$
  A\subseteq\mathcal{S}
  \longmapsto
  N_{A}
  =
  \bigcup_{S\in A}2^{\Omega\setminus S}
$$
has image exactly the set of subcomplexes of the tetrahedron boundary
$\partial\Delta^{3}$, with both the void complex and the complex
consisting only of the empty face included.
\end{theorem}
\begin{proof}
For every nonempty support $S\subseteq\Omega$, the complement
$\Omega\setminus S$ is a proper face of $\Delta^{3}$. Hence $N_{A}$ is
the downward closure of a family of proper faces, and therefore is a
subcomplex of $\partial\Delta^{3}$.

Conversely, let $K\subseteq\partial\Delta^{3}$ be a subcomplex. For
each maximal face $F$ of $K$, set $S_{F}=\Omega\setminus F$. These
supports are nonempty. If $A=\{S_{F}:F\in\operatorname{Facets}(K)\}$,
then
$$
  N_{A}
  =
  \bigcup_{F\in\operatorname{Facets}(K)}2^{F}
  =
  K.
$$
Thus every subcomplex of $\partial\Delta^{3}$ is realized.
\end{proof}

\begin{theorem}[Support-family fiber]
\label{thm:philosophy-support-family-fiber}
Let $K\subseteq\partial\Delta^{3}$ be a subcomplex, and let $|K|$
denote its number of faces, including the empty face. The generating
function for support families $A$ with $N_{A}=K$ is
$$
  G_{K}(z)
  =
  z^{|\operatorname{Facets}(K)|}
  (1+z)^{|K|-|\operatorname{Facets}(K)|}.
$$
\end{theorem}
\begin{proof}
Write
$$
  C_{A}=\{\Omega\setminus S:S\in A\}.
$$
Then $N_{A}$ is the downward closure of $C_{A}$. The equality $N_{A}=K$
holds exactly when $C_{A}\subseteq K$ and $C_{A}$ contains every maximal
face of $K$. The maximal faces are forced, and the remaining
$|K|-|\operatorname{Facets}(K)|$ faces may be chosen independently.
\end{proof}

\begin{corollary}[Subcomplex-weighted support polynomials]
\label{cor:philosophy-subcomplex-weighted-support-polynomials}
For $q=-1,0,1,2$,
$$
  Q_{q}(z)
  =
  \sum_{K\subseteq\partial\Delta^{3}}
  \dim_{k}\widetilde{H}_{q}(K;k)
  z^{|\operatorname{Facets}(K)|}
  (1+z)^{|K|-|\operatorname{Facets}(K)|}.
$$
\end{corollary}
\begin{proof}
Group the defining sum for $Q_{q}(z)$ by the realized subcomplex
$K=N_{A}$ and apply \autoref{thm:philosophy-support-family-fiber}.
\end{proof}

\begin{theorem}[Tetrahedron-boundary subcomplex census]
\label{thm:philosophy-tetrahedron-boundary-subcomplex-census}
The boundary $\partial\Delta^{3}$ has $167$ subcomplexes under the
convention of \autoref{thm:philosophy-support-to-nerve-realization}.
Their reduced Betti signatures are:
$$
\begin{array}{c|r}
  \text{signature} & \text{subcomplex count}\\
  \hline
  (1,0,0,0) & 2\\
  (0,0,0,0) & 64\\
  (0,1,0,0) & 37\\
  (0,2,0,0) & 10\\
  (0,3,0,0) & 1\\
  (0,0,1,0) & 37\\
  (0,0,2,0) & 10\\
  (0,0,3,0) & 1\\
  (0,1,1,0) & 4\\
  (0,0,0,1) & 1.
\end{array}
$$
Moreover,
$$
  \sum_{K\subseteq\partial\Delta^{3}}
  2^{|K|-|\operatorname{Facets}(K)|}
  =
  32768.
$$
\end{theorem}
\begin{proof}
Enumerate the downward-closed families of proper faces of a four-vertex
simplex and compute reduced homology for each subcomplex. The displayed
rows sum to $167$. The final identity is the specialization
$G_{K}(1)$ summed over all realized subcomplexes; by
\autoref{thm:philosophy-support-family-fiber}, it counts all
$2^{15}$ support families.
\end{proof}

\begin{theorem}[Support-to-nerve Betti formula]
\label{thm:philosophy-support-to-nerve-betti-formula}
The complete Hochster diagonal series of $k[\mathcal{B}]$ satisfy
$$
  B_{q}(x)
  =
  \sum_{K\subseteq\partial\Delta^{3}}
  \dim_{k}\widetilde{H}_{q}(K;k)
  (2x+x^{2})^{|\operatorname{Facets}(K)|}
  (1+2x+x^{2})^{|K|-|\operatorname{Facets}(K)|}.
$$
Equivalently,
$$
  B_{q}(x)
  =
  \sum_{K\subseteq\partial\Delta^{3}}
  \dim_{k}\widetilde{H}_{q}(K;k)
  x^{|\operatorname{Facets}(K)|}
  (x+2)^{|\operatorname{Facets}(K)|}
  (1+x)^{2(|K|-|\operatorname{Facets}(K)|)}.
$$
\end{theorem}
\begin{proof}
Combine
\autoref{cor:philosophy-subcomplex-weighted-support-polynomials} with
the codon lift substitution of
\autoref{thm:philosophy-support-lift-betti-transform}. The two displayed
forms are identical because $2x+x^{2}=x(x+2)$ and
$1+2x+x^{2}=(1+x)^{2}$.
\end{proof}

\begin{corollary}[Top diagonal from the boundary complex]
\label{cor:philosophy-top-diagonal-from-boundary-complex}
The top Hochster diagonal is contributed by the single subcomplex
$K=\partial\Delta^{3}$. Since this complex has four facets and fifteen
proper faces,
$$
  B_{2}(x)
  =
  (2x+x^{2})^{4}(1+2x+x^{2})^{11}
  =
  x^{4}(x+2)^{4}(1+x)^{22}.
$$
In particular,
$$
  \beta_{27,30}=1.
$$
\end{corollary}
\begin{proof}
Among subcomplexes of $\partial\Delta^{3}$, nonzero
$\widetilde{H}_{2}$ occurs only for $\partial\Delta^{3}$ itself. The
formula is the $q=2$ specialization of
\autoref{thm:philosophy-support-to-nerve-betti-formula}.
\end{proof}
